\section{Phase Stiffness, Quasi-Two-Dimensional Coherence, and Operational Proxies}
  \label{sec:phase-stiffness}

  \subsection{Phase Stiffness as Relaxation Capacity}
    \label{subsec:stiffness-relaxation}

    In the strongly correlated regime,
    the decisive quantity controlling the superconducting transition
    is not the amplitude gap
    $\Delta_\Pi^{\mathrm{amplitude}}$,
    but the phase stiffness $\rho_s$.

    Within the projective framework,
    $\rho_s$ measures the relaxation capacity of the projected substrate,
    that is,
    its ability to redistribute local fiber phase mismatches
    without generating topological defects.
    We therefore identify $\rho_s$
    as the operational proxy for the relaxation flux
    \begin{equation}
      \mathcal{R}
      \equiv
      \rho_s .
      \label{eq:relaxation-proxy}
    \end{equation}

    In doped Mott systems,
    the density of active composites
    is proportional to the carrier concentration $\delta$.
    Combined with Eq.~\ref{eq:gap-frustration},
    this yields to leading order
    \begin{equation}
      \rho_s^{(0)}
      \propto
      \delta J .
      \label{eq:rho-deltaJ}
    \end{equation}

    Equation~\ref{eq:rho-deltaJ}
    provides the structural origin
    of the experimentally observed Uemura scaling
    in underdoped cuprates.

  \subsection{Quasi-Two-Dimensional Phase Transition}
    \label{subsec:quasi-2d}

    Cuprates and related materials
    are strongly anisotropic layered systems.
    The dominant projective substrate
    is provided by the CuO$_2$ planes,
    while the interlayer coupling $J_\perp$
    is significantly smaller than the in-plane exchange scale $J$.

    In this quasi-two-dimensional limit,
    phase coherence is established within each layer
    through vortex unbinding.
    The superconducting transition
    is therefore governed by the
    Berezinskii--Kosterlitz--Thouless (BKT) criterion,
    \begin{equation}
      k_B T_c
      \simeq
      \frac{\pi}{2}
      \rho_s(T_c^-),
      \label{eq:BKT}
    \end{equation}
    where $\rho_s(T_c^-)$ denotes the renormalized stiffness
    evaluated at the vortex scale.

    The system is not strictly two-dimensional.
    Finite interlayer coupling
    converts the BKT transition
    into a true three-dimensional transition
    at a slightly higher temperature.
    The quasi-2D description applies
    in the strongly anisotropic regime
    \begin{equation}
      J_\perp
      \ll
      k_B T_c
      \ll
      J ,
    \end{equation}
    which is satisfied in cuprate families
    over a broad doping range.

  \subsection{Frustration Density as Observable Proxy}
    \label{subsec:frustration-observable}

    To avoid purely definitional constructs,
    the frustration density $\mathcal{F}$
    must be independently accessible from experiment.

    We use as first-order proxy
    \begin{equation}
      \mathcal{F}
      \propto
      \frac{S(\pi,\pi)}{S_{\max}},
      \label{eq:F-proxy}
    \end{equation}
    where $S(\pi,\pi)$ is the magnetic structure factor
    measured by neutron scattering or RIXS,
    and $S_{\max}$ is the fully ordered reference value.

    This proportionality is a leading-order assumption.
    Non-linear corrections would modify numerical prefactors
    but do not alter the qualitative phase selection mechanism
    unless fine-tuned cancellations occur.

  \subsection{Three Regimes from the Ratio \texorpdfstring{$\mathcal{F}/\mathcal{R}$}{F/R}}
    \label{subsec:F-over-R}

    The interplay between frustration density $\mathcal{F}$
    and relaxation capacity $\mathcal{R}=\rho_s$
    defines three qualitatively distinct regimes.

    \begin{enumerate}
      \item
      \textbf{Antiferromagnetic regime:}
      \(
      \mathcal{F} \rightarrow 1,
      \quad
      \rho_s \rightarrow 0
      \).
      Long-range AF order dominates,
      and composites cannot delocalize coherently.

      \item
      \textbf{Superconducting regime:}
      Intermediate $\mathcal{F}$
      with finite $\rho_s$.
      Frustration is sufficient to stabilize composites,
      while relaxation capacity allows global phase locking.

      \item
      \textbf{Glass-like or incoherent metallic regime:}
      \(
      \mathcal{F} \rightarrow 0,
      \quad
      \rho_s \rightarrow 0
      \)
      with strong disorder.
      Frustration is too weak or too fragmented
      to stabilize composites coherently.
    \end{enumerate}

    These regimes are experimentally distinguishable.
    A material exhibiting intermediate $S(\pi,\pi)$
    but vanishing $\rho_s$
    must possess an additional mechanism suppressing
    phase coherence,
    such as strong disorder or orbital mixing.
    Conversely,
    superconductivity in a system
    with negligible $(\pi,\pi)$ frustration
    requires a different frustration channel,
    providing a falsifiable diagnostic criterion.

  \subsection{Emergent Dome Structure}
    \label{subsec:dome-final}

    The superconducting dome arises naturally
    from the doping dependence of both $\mathcal{F}$ and $\rho_s$.

    At low doping,
    $\mathcal{F}$ is large but $\rho_s$ is small.
    At optimal doping,
    $\mathcal{F}$ is reduced
    while $\rho_s \propto \delta J$ is maximal.
    At overdoping,
    $\mathcal{F}$ becomes too small
    to sustain a significant projective gap,
    and $T_c$ decreases accordingly.

    No additional phenomenological tuning parameter
    is required to generate the dome structure.
    It emerges from the coupled evolution
    of frustration amplitude and phase stiffness.
